\changecolours{ScoutGreen}
\section[Especialização e Generalização]{Especialização, Generalização e a Notação EER}

% ------------------------------------------------------------------------------
\twocolframe[title={Abordagens: Top-Down vs. Bottom-Up},leftcol=7.1cm,rightcol=6.9cm]{%
  \alert{Especialização (Abordagem Top-Down)}\par
  \itemseps[topsepi=0.5mm,itemsepi=1mm,labelsep=2mm,leftmargini=3mm]
  \begin{itemize}
    \item \textbf{Processo:} Parte-se de um tipo de entidade genérico existente e criam-se subclasses baseadas em características distintivas.
    \item \textbf{Quando Ocorre?}
      Quando o analista percebe que certos subgrupos de instâncias possuem propriedades específicas ou regras de negócio exclusivas.
    \item \textbf{Exemplo:} \texttt{FUNCIONARIO} especializa-se em \texttt{ENGENHEIRO}, \texttt{SECRETARIO} e \texttt{TECNICO}.
  \end{itemize}
}{%
  \alert{Generalização (Abordagem Bottom-Up)}\par
  \itemseps[topsepi=0.5mm,itemsepi=1mm,labelsep=2mm,leftmargini=3mm]
  \begin{itemize}
    \item \textbf{Processo:} Identificam-se múltiplos tipos de entidades já existentes com características comuns e cria-se uma superclasse genérica para unificá-los.
    \item \textbf{Quando Ocorre?}
      Em fases de consolidação de requisitos onde modelos parciais são integrados.
    \item \textbf{Exemplo:} Unificar \texttt{CARRO}, \texttt{CAMINHAO} e \texttt{MOTO} na superclasse \texttt{VEICULO}.
  \end{itemize}
}

% ------------------------------------------------------------------------------
\begin{frame}{Representação Gráfica Canônica no Diagrama EER}
  \vspace{0.1cm}
  \begin{center}
    \tikz[scale=0.75, transform shape]{
      % Superclasse
      \node[draw=ScoutGreen, fill=ScoutGreen!15, thick, font=\bfseries\sffamily\scriptsize, minimum width=3.0cm, minimum height=0.65cm] (super) {VEICULO};
      \node[draw=ScoutNavy, fill=ScoutNavy!10, ellipse, font=\tiny\sffamily, above left=0.4cm and 0.3cm of super] (pk) {\underline{placa}};
      \node[draw=ScoutNavy, fill=ScoutNavy!10, ellipse, font=\tiny\sffamily, above right=0.4cm and 0.3cm of super] (fab) {marca};
      \draw[thick, ScoutNavy] (super) -- (pk);
      \draw[thick, ScoutNavy] (super) -- (fab);

      % Circulo EER (Letra d indicando disjuncao)
      \node[draw=ScoutNavy, fill=ScoutNavy!20, circle, thick, inner sep=2.5pt, font=\bfseries\sffamily\scriptsize, below=0.65cm of super] (circ) {d};

      % Subclasses
      \node[draw=ScoutGreen, fill=ScoutGreen!15, thick, font=\bfseries\sffamily\scriptsize, minimum width=2.2cm, minimum height=0.65cm, below left=0.7cm and 1.0cm of circ] (sub1) {CARRO};
      \node[draw=ScoutNavy, fill=ScoutNavy!10, ellipse, font=\tiny\sffamily, below=0.35cm of sub1] (atr_sub1) {qtd\_portas};
      \draw[thick, ScoutNavy] (sub1) -- (atr_sub1);

      \node[draw=ScoutGreen, fill=ScoutGreen!15, thick, font=\bfseries\sffamily\scriptsize, minimum width=2.2cm, minimum height=0.65cm, below right=0.7cm and 1.0cm of circ] (sub2) {CAMINHAO};
      \node[draw=ScoutNavy, fill=ScoutNavy!10, ellipse, font=\tiny\sffamily, below=0.35cm of sub2] (atr_sub2) {capacidade\_carga};
      \draw[thick, ScoutNavy] (sub2) -- (atr_sub2);

      % Linhas conectando superclasse ao circulo e circulo as subclasses
      \draw[thick, double, double distance=1.5pt, black!80] (super) -- (circ);
      \draw[thick, black!80] (circ) -| (sub1);
      \draw[thick, black!80] (circ) -| (sub2);

      % Simbolo de subconjunto U na ligacao
      \node[font=\footnotesize, left=0.2cm of sub1.north, yshift=0.3cm] {$\cup$};
      \node[font=\footnotesize, right=0.2cm of sub2.north, yshift=0.3cm] {$\cup$};
    }
  \end{center}
  \vspace{0.05cm}
  \centering{\tiny \textbf{Notação EER:} A linha dupla para o círculo indica especialização total; a letra ``d'' indica disjunção exclusiva.}
\end{frame}
